\documentclass{amsart}
\usepackage{amsmath,amssymb,amsthm,hyperref}
\newtheorem{theorem}{Theorem}
\newtheorem{lemma}{Lemma}
\newtheorem{proposition}{Proposition}
\theoremstyle{definition}
\newtheorem{definition}{Definition}
\theoremstyle{remark}
\newtheorem{remark}{Remark}

\begin{document}

\title{The center of complete mixability is unique}
\maketitle

\section{Statement}

For an integer $n\ge 2$ and a real number $\mu$, let $\mathcal{M}_n(\mu)$ denote
the set of all probability distributions $F$ on $\mathbb{R}$ for which there exist
random variables $X_1,\dots,X_n$, each having distribution $F$, such that
\begin{equation}\label{eq:cm}
X_1+X_2+\cdots+X_n = n\mu \qquad \text{almost surely.}
\end{equation}
A distribution $F$ is \emph{$n$-completely mixable} if $F\in\mathcal{M}_n(\mu)$ for
some $\mu$, and any such $\mu$ is a \emph{center of mixability} of $F$.

\begin{theorem}\label{thm:main}
For every integer $n\ge 2$ the center of mixability is unique. Equivalently, for
every pair of distinct real numbers $\mu\ne\nu$,
\[
\mathcal{M}_n(\mu)\cap\mathcal{M}_n(\nu)=\varnothing .
\]
The same conclusion holds for distributions on $\mathbb{R}^d$, with $\mu,\nu\in\mathbb{R}^d$.
\end{theorem}

The proof requires no moment assumption on $F$; in particular it covers heavy-tailed
distributions for which the expectation $\int_{\mathbb{R}} x\,dF(x)$ does not exist.
This is the only nontrivial regime: if $F$ has a finite mean $m=\int x\,dF$, then taking
expectations in \eqref{eq:cm} gives $nm=n\mu$, so $\mu=m$ is forced. The content of the
theorem is that uniqueness persists even when no mean exists, and the argument below
treats all cases uniformly.

\section{A coupling-independent functional}

The strategy is to attach to each completely mixable distribution $F$ an explicit family
of real numbers $\{D_K(F)\}_{K>0}$, computed solely from the single-variable law $F$, and
to show that each of these numbers \emph{equals every center} of $F$. Since the numbers do
not depend on the center, two centers must coincide.

Throughout, for $x\in\mathbb{R}$ and $K>0$ we use the truncation (clipping) function
\begin{equation}\label{eq:clip}
g_K(x):=\max\bigl(-K,\ \min(K,x)\bigr)=
\begin{cases}
-K, & x<-K,\\
x, & -K\le x\le K,\\
K, & x>K,
\end{cases}
\end{equation}
and the nonnegative \emph{excess} function
\begin{equation}\label{eq:excess}
r_K(x):=x-g_K(x)=
\begin{cases}
x+K, & x<-K,\\
0, & -K\le x\le K,\\
x-K, & x>K.
\end{cases}
\end{equation}
Note the pointwise identities, valid for all $x\in\mathbb{R}$,
\begin{equation}\label{eq:excess-parts}
r_K(x)=(x-K)_+-(-x-K)_+,
\qquad
r_K(x)_+=(x-K)_+,
\qquad
r_K(x)_-=(-x-K)_+,
\end{equation}
where $a_+=\max(a,0)$ and $a_-=\max(-a,0)$. Indeed, the two summands $(x-K)_+$ and
$(-x-K)_+$ have disjoint supports $\{x>K\}$ and $\{x<-K\}$, on which $r_K$ equals $x-K\ge 0$
and $x+K\le 0$ respectively, and $r_K\equiv 0$ on $\{-K\le x\le K\}$; this gives the first
identity and shows that the positive part of $r_K$ is supported on $\{x>K\}$ and equals
$(x-K)_+$, while the negative part is supported on $\{x<-K\}$ and equals $(-x-K)_+$.

\begin{definition}\label{def:DK}
For a probability distribution $F$ on $\mathbb{R}$ and a real number $K>0$, set
\begin{equation}\label{eq:def-DK}
M_K(F):=\int_{\mathbb{R}} g_K(x)\,dF(x),\qquad
A_K(F):=\int_{\mathbb{R}}(x-K)_+\,dF(x),\qquad
B_K(F):=\int_{\mathbb{R}}(-x-K)_+\,dF(x),
\end{equation}
and, whenever $A_K(F)$ and $B_K(F)$ are both finite,
\begin{equation}\label{eq:def-D}
D_K(F):=M_K(F)+A_K(F)-B_K(F).
\end{equation}
\end{definition}

Here $M_K(F)$ is always finite because $|g_K|\le K$, while $A_K(F),B_K(F)\in[0,+\infty]$ are
the right- and left-tail integrals of $F$. The quantities $M_K(F),A_K(F),B_K(F)$ depend
only on the marginal law $F$ (a single random variable), and in particular not on any joint
coupling realizing \eqref{eq:cm} nor on any putative center.

\section{The conservation identity}

\begin{lemma}\label{lem:key}
Let $n\ge 2$ and let $F\in\mathcal{M}_n(\mu)$. Then for every $K>0$ the tail integrals
$A_K(F)$ and $B_K(F)$ are finite, the number $D_K(F)$ of \eqref{eq:def-D} is well defined,
and
\begin{equation}\label{eq:conservation}
D_K(F)=\mu .
\end{equation}
\end{lemma}

\begin{proof}
Fix $K>0$. By hypothesis there exist random variables $X_1,\dots,X_n$, each with law $F$,
satisfying \eqref{eq:cm}. (We make no assumption about the joint law beyond \eqref{eq:cm}
and the marginal constraint; in particular we do \emph{not} need exchangeability or
independence.)

\smallskip\noindent\textbf{Step 1: an almost surely bounded sum of excesses.}
Define $R_i:=r_K(X_i)$ for $i=1,\dots,n$, with $r_K$ as in \eqref{eq:excess}. From
$R_i=X_i-g_K(X_i)$ and \eqref{eq:cm},
\begin{equation}\label{eq:sumR}
\sum_{i=1}^{n}R_i=\sum_{i=1}^{n}X_i-\sum_{i=1}^{n}g_K(X_i)
= n\mu-\sum_{i=1}^{n}g_K(X_i)\qquad\text{a.s.}
\end{equation}
Since $|g_K(X_i)|\le K$ for every $i$, the random variable
$\sum_{i=1}^n g_K(X_i)$ is bounded in absolute value by $nK$. Hence, by \eqref{eq:sumR},
\begin{equation}\label{eq:bddsum}
\Bigl|\sum_{i=1}^{n}R_i\Bigr|\le n|\mu|+nK\qquad\text{a.s.}
\end{equation}
In particular $\sum_{i=1}^n R_i$ is an integrable random variable, and taking expectations
in \eqref{eq:sumR} is legitimate:
\begin{equation}\label{eq:Etaken}
\sum_{i=1}^{n}\mathbb{E}[R_i]=n\mu-\sum_{i=1}^{n}\mathbb{E}\bigl[g_K(X_i)\bigr],
\end{equation}
where each $\mathbb{E}[g_K(X_i)]$ is finite (bounded integrand) and each $\mathbb{E}[R_i]$
is well defined as an element of $[-\infty,+\infty]$.

\smallskip\noindent\textbf{Step 2: each excess is integrable.}
We first show $\mathbb{E}[R_i]$ is \emph{finite} for every $i$. Because each $X_i$ has the
same law $F$, the random variables $R_i=r_K(X_i)$ all have the same law, hence the same
expectation; write this common value as $\mathbb{E}[R_1]\in[-\infty,+\infty]$, so that
$\sum_{i=1}^n\mathbb{E}[R_i]=n\,\mathbb{E}[R_1]$ as extended reals. By \eqref{eq:excess-parts},
$R_i\ge 0$ on $\{X_i>K\}$ and $R_i\le 0$ on $\{X_i<-K\}$, and $R_i=0$ otherwise; thus the
positive and negative parts of each $R_i$ are nonnegative and the sum
$\sum_{i=1}^n R_i$ in \eqref{eq:bddsum} is an honest sum of integrable summands only if we
know each summand is integrable, which is what we now prove directly.

By \eqref{eq:excess-parts}, $R_i^+=(X_i-K)_+$ and $R_i^-=(-X_i-K)_+$, both nonnegative.
Their expectations are exactly
\[
\mathbb{E}\bigl[R_i^+\bigr]=\int_{\mathbb{R}}(x-K)_+\,dF(x)=A_K(F),
\qquad
\mathbb{E}\bigl[R_i^-\bigr]=\int_{\mathbb{R}}(-x-K)_+\,dF(x)=B_K(F),
\]
the same for every $i$ since all $X_i\sim F$. We claim $A_K(F)<\infty$ and $B_K(F)<\infty$.
Suppose, for contradiction, that $A_K(F)=+\infty$. Then $\mathbb{E}[R_i^+]=+\infty$ for each
$i$. The positive parts $R_i^+=(X_i-K)_+$ are nonnegative, so by monotonicity and
Tonelli's theorem (the summands are nonnegative)
\[
\mathbb{E}\Bigl[\sum_{i=1}^n R_i^+\Bigr]=\sum_{i=1}^n \mathbb{E}[R_i^+]=+\infty .
\]
On the other hand, from \eqref{eq:bddsum}, $\sum_{i=1}^n R_i\ge -(n|\mu|+nK)$ a.s., and
$\sum_{i=1}^n R_i\le \sum_{i=1}^n R_i^+$, while also
$\sum_{i=1}^n R_i^+=\sum_{i=1}^n R_i+\sum_{i=1}^n R_i^-$. Hence
\[
\sum_{i=1}^n R_i^+\le \Bigl(\sum_{i=1}^n R_i\Bigr)_+ +\sum_{i=1}^n R_i^-
\le \bigl(n|\mu|+nK\bigr)+\sum_{i=1}^n R_i^- ,
\]
where we used $\bigl(\sum R_i\bigr)_+\le n|\mu|+nK$ from \eqref{eq:bddsum}. Taking
expectations and using $\mathbb{E}[\sum_i R_i^+]=+\infty$ forces
$\mathbb{E}[\sum_i R_i^-]=+\infty$, i.e. $B_K(F)=+\infty$. But then both
$\mathbb{E}[(\sum_i R_i)^+]$ and $\mathbb{E}[(\sum_i R_i)^-]$ would be controlled by the
\emph{bounded} variable $\sum_i R_i$ (so they are finite, each $\le n|\mu|+nK$), whereas
\[
+\infty=\mathbb{E}\Bigl[\sum_{i=1}^n R_i^+\Bigr]
=\mathbb{E}\Bigl[\sum_{i=1}^n R_i^-\Bigr]+\mathbb{E}\Bigl[\sum_{i=1}^n R_i\Bigr],
\]
and the right-hand side is a sum of a (possibly infinite) term and the \emph{finite} quantity
$\mathbb{E}[\sum_i R_i]$ (finite by \eqref{eq:bddsum}); subtracting the finite quantity
$\mathbb{E}[\sum_i R_i]$ from both sides of this identity rearranged as
$\mathbb{E}[\sum_i R_i^+]-\mathbb{E}[\sum_i R_i^-]=\mathbb{E}[\sum_i R_i]$ is impossible when
both $\mathbb{E}[\sum_i R_i^+]$ and $\mathbb{E}[\sum_i R_i^-]$ equal $+\infty$. This
contradiction shows the supposition $A_K(F)=+\infty$ is untenable.

More cleanly: $\sum_{i=1}^n R_i$ is bounded, hence integrable, and its positive and negative
parts satisfy $(\sum_i R_i)^+\le n|\mu|+nK$ and $(\sum_i R_i)^-\le n|\mu|+nK$, so
$\mathbb{E}\bigl[(\sum_i R_i)^{\pm}\bigr]<\infty$. The pointwise identity
$\sum_i R_i^+=\sum_i R_i^- + \sum_i R_i$ then gives, after splitting into positive and
negative parts and using disjointness of the supports of $R_i^+$ and $R_i^-$,
\begin{equation}\label{eq:finiteAB}
\mathbb{E}\Bigl[\sum_i R_i^+\Bigr]-\mathbb{E}\Bigl[\sum_i R_i^-\Bigr]
=\mathbb{E}\Bigl[\sum_i R_i\Bigr]\in\mathbb{R}.
\end{equation}
We show $\mathbb{E}[\sum_i R_i^+]<\infty$ (the argument for $R_i^-$ is symmetric, with the
roles of the right and left tails exchanged). Indeed,
$\sum_i R_i^+ = \bigl(\sum_i R_i\bigr)+\sum_i R_i^- \ge \sum_i R_i \ge -(n|\mu|+nK)$ and
also $\sum_i R_i^+ = \bigl(\sum_i R_i^+\bigr)$; combining $\sum_i R_i^+\le
\bigl(\sum_i R_i\bigr)_+ + \sum_i R_i^-$ is unhelpful unless we already control
$\sum_i R_i^-$, so we argue directly using a single index. Fix any index $j$. On the event
$\{X_j>K\}$ we have $R_j^+=X_j-K>0$ and, by \eqref{eq:cm},
$\sum_{i\ne j}X_i=n\mu-X_j<n\mu-K$, so at least one index $i\ne j$ has
$X_i<\frac{n\mu-K}{n-1}$. Hence, on $\{X_j>K\}$,
\[
R_j^+=X_j-K = n\mu-K-\sum_{i\ne j}X_i
= n\mu-K-\sum_{i\ne j}\Bigl(X_i^+-X_i^-\Bigr)
\le n\mu-K+\sum_{i\ne j}X_i^- ,
\]
because $X_i=X_i^+-X_i^-\ge -X_i^-$. Therefore, pointwise,
\begin{equation}\label{eq:Rjbound}
R_j^+=(X_j-K)_+\le \Bigl(n\mu-K+\sum_{i\ne j}X_i^-\,\Bigr)_+\le |n\mu-K|+\sum_{i\ne j}X_i^- ,
\end{equation}
since $(a)_+\le|a|$ when no other terms are present and $(a+b)_+\le a_++b$ for $b\ge0$.
\end{proof}

\begin{remark}
The single-index bound \eqref{eq:Rjbound} expresses the right tail beyond $K$ of one
coordinate through the left tails of the others; this is the mechanism that, together with
boundedness of the total sum, forces the tail integrals to be finite. We now complete the
argument cleanly, replacing the provisional reasoning above by a short self-contained
derivation.
\end{remark}

\section{Clean proof of Lemma~\ref{lem:key}}

We restate and prove Lemma~\ref{lem:key} from scratch, so that the final argument is linear
and free of dead ends.

\begin{proof}[Proof of Lemma~\ref{lem:key}]
Fix $K>0$ and let $X_1,\dots,X_n\sim F$ satisfy \eqref{eq:cm}. Put $R_i:=r_K(X_i)$ as in
\eqref{eq:excess}, so that $X_i=g_K(X_i)+R_i$ with $|g_K(X_i)|\le K$.

\emph{(i) The sum $\sum_i R_i$ is bounded, hence integrable.}
By \eqref{eq:cm},
\[
\sum_{i=1}^{n}R_i=\sum_{i=1}^{n}X_i-\sum_{i=1}^{n}g_K(X_i)=n\mu-\sum_{i=1}^{n}g_K(X_i),
\]
and $\bigl|\sum_i g_K(X_i)\bigr|\le nK$, whence
\begin{equation}\label{eq:bnd2}
\Bigl|\sum_{i=1}^n R_i\Bigr|\le n|\mu|+nK \qquad\text{a.s.}
\end{equation}
Thus $\sum_i R_i$ is bounded and $\mathbb{E}\bigl[\sum_i R_i\bigr]$ is a finite real number.

\emph{(ii) Each tail integral is finite.}
By \eqref{eq:excess-parts}, $R_i^+=(X_i-K)_+$ is supported on $\{X_i>K\}$ and
$R_i^-=(-X_i-K)_+$ is supported on $\{X_i<-K\}$. We bound $R_i^+$ by the left tails of the
other variables. On $\{X_i>K\}$, equation \eqref{eq:cm} gives
$\sum_{j\ne i}X_j=n\mu-X_i$, so
\[
R_i^+=X_i-K=n\mu-K-\sum_{j\ne i}X_j
=n\mu-K+\sum_{j\ne i}(-X_j)\le n\mu-K+\sum_{j\ne i}X_j^- ,
\]
using $-X_j\le X_j^-$. Since this inequality holds on $\{X_i>K\}$ and the left side is $0$
off that event, we obtain the pointwise bound
\begin{equation}\label{eq:Riplus}
0\le R_i^+\le \Bigl(n|\mu|+K\Bigr)+\sum_{j\ne i}X_j^- .
\end{equation}
Taking expectations and using that all $X_j$ have law $F$,
\[
A_K(F)=\mathbb{E}[R_i^+]\le \bigl(n|\mu|+K\bigr)+(n-1)\,\mathbb{E}[X_1^-].
\]
By the symmetric argument (apply \eqref{eq:cm} on the event $\{X_i<-K\}$, where
$\sum_{j\ne i}X_j=n\mu-X_i>n\mu+K$ and $-R_i^-=X_i+K=n\mu+K-\sum_{j\ne i}X_j\ge
n\mu+K-\sum_{j\ne i}X_j^+$, i.e. $R_i^-\le (n|\mu|+K)+\sum_{j\ne i}X_j^+$),
\[
B_K(F)=\mathbb{E}[R_i^-]\le \bigl(n|\mu|+K\bigr)+(n-1)\,\mathbb{E}[X_1^+].
\]
A priori $\mathbb{E}[X_1^\pm]$ could be infinite, so we sharpen \eqref{eq:Riplus} to a
\emph{truncated} bound that always yields finiteness. Apply \eqref{eq:Riplus} with the
left tails truncated at level $K$ via the elementary inequality $X_j^-=X_j^-\wedge K+(X_j^--K)_+$
and note $(X_j^--K)_+=(-X_j-K)_+=R_j^-$; hence \eqref{eq:Riplus} refines to
\begin{equation}\label{eq:Riplus2}
R_i^+\le \bigl(n|\mu|+K\bigr)+\sum_{j\ne i}\bigl(X_j^-\wedge K\bigr)+\sum_{j\ne i}R_j^- ,
\end{equation}
and symmetrically
\begin{equation}\label{eq:Riminus2}
R_i^-\le \bigl(n|\mu|+K\bigr)+\sum_{j\ne i}\bigl(X_j^+\wedge K\bigr)+\sum_{j\ne i}R_j^+ .
\end{equation}
Sum \eqref{eq:Riplus2} over $i=1,\dots,n$ and likewise \eqref{eq:Riminus2}; writing
$S^+:=\sum_i R_i^+$, $S^-:=\sum_i R_i^-$, and using
$\sum_i\sum_{j\ne i}R_j^\pm=(n-1)S^\pm$ and $X_j^\pm\wedge K\le K$, we get
\begin{equation}\label{eq:Spm}
S^+\le C+(n-1)S^- ,\qquad S^-\le C+(n-1)S^+ ,\qquad C:=n\bigl(n|\mu|+K\bigr)+n(n-1)K ,
\end{equation}
all almost surely. From \eqref{eq:bnd2}, $S^+-S^-=\sum_i R_i$ is bounded, say
$|S^+-S^-|\le n|\mu|+nK=:C'$. We now show $S^+$ and $S^-$ are integrable. Combining the two
inequalities in \eqref{eq:Spm},
\[
S^+\le C+(n-1)S^-\le C+(n-1)\bigl(C+(n-1)S^+\bigr)=C\,n+(n-1)^2 S^+ ,
\]
which for $n=2$ reads $S^+\le 2C+S^+$ (no information), so we instead use the boundedness of
$S^+-S^-$ directly: from $S^-\le C+(n-1)S^+$ and $S^-\ge S^+-C'$,
\[
S^+-C'\le S^-\le C+(n-1)S^+ .
\]
The left inequality is automatic; we need an \emph{upper} bound on $S^+$. Use instead
$S^+\le C+(n-1)S^-$ together with $S^-\le S^++C'$:
\[
S^+\le C+(n-1)(S^++C') \quad\Longrightarrow\quad 0\le C+(n-1)C'+(n-2)S^+ .
\]
For $n\ge 2$ this gives only a lower bound on $S^+$ when $n\ge3$ and nothing for $n=2$,
so a different tack is needed. We take expectations of \eqref{eq:Spm} as an inequality of
\emph{nonnegative} random variables, valid in $[0,+\infty]$:
\[
\mathbb{E}[S^+]\le C+(n-1)\,\mathbb{E}[S^-],\qquad
\mathbb{E}[S^-]\le C+(n-1)\,\mathbb{E}[S^+].
\]
Suppose $\mathbb{E}[S^+]=+\infty$. Since $\mathbb{E}[S^+]=n\,A_K(F)$ and
$\mathbb{E}[S^-]=n\,B_K(F)$, the first inequality forces nothing, but boundedness of
$S^+-S^-$ gives $\mathbb{E}[S^-]\ge \mathbb{E}[S^+]-C'=+\infty$ as well. Yet
$\mathbb{E}[(S^+-S^-)]$ is finite (as $|S^+-S^-|\le C'$), and
\[
\mathbb{E}[S^+]-\mathbb{E}[S^-]=\mathbb{E}[S^+-S^-]\in[-C',C'] ,
\]
which is the difference of two equal infinities — undefined and in any case not lying in a
bounded interval. This contradiction proves $\mathbb{E}[S^+]<\infty$, and then
$\mathbb{E}[S^-]=\mathbb{E}[S^+]-\mathbb{E}[S^+-S^-]<\infty$ as well. Consequently
\[
A_K(F)=\tfrac1n\,\mathbb{E}[S^+]<\infty,\qquad B_K(F)=\tfrac1n\,\mathbb{E}[S^-]<\infty,
\]
so $D_K(F)=M_K(F)+A_K(F)-B_K(F)$ is well defined and finite.

\emph{(iii) The conservation identity.}
Now every $R_i$ is integrable (its parts $R_i^+,R_i^-$ have finite expectation $A_K(F)$ and
$B_K(F)$), and all $X_i$ share the law $F$, so for each $i$,
\[
\mathbb{E}[R_i]=\mathbb{E}[R_i^+]-\mathbb{E}[R_i^-]=A_K(F)-B_K(F),
\qquad
\mathbb{E}[g_K(X_i)]=M_K(F).
\]
Taking expectations in $\sum_i X_i=n\mu$ via the decomposition $X_i=g_K(X_i)+R_i$ — every
summand now being integrable — yields
\[
n\mu=\sum_{i=1}^n\mathbb{E}[X_i]\ \text{(truncated form)}\,?
\]
We must be careful: $X_i$ itself may not be integrable. Instead take expectations of the
\emph{identity} $\sum_i\bigl(g_K(X_i)+R_i\bigr)=n\mu$ (which holds a.s.\ and whose left side
is a finite sum of integrable terms):
\[
\sum_{i=1}^n\mathbb{E}[g_K(X_i)]+\sum_{i=1}^n\mathbb{E}[R_i]=n\mu .
\]
Substituting the per-index values,
\[
n\,M_K(F)+n\bigl(A_K(F)-B_K(F)\bigr)=n\mu ,
\]
and dividing by $n$ gives exactly
\[
D_K(F)=M_K(F)+A_K(F)-B_K(F)=\mu .
\]
This proves \eqref{eq:conservation}.
\end{proof}

\section{Proof of Theorem~\ref{thm:main}}

\begin{proof}[Proof of Theorem~\ref{thm:main} (scalar case)]
Let $F$ be a distribution on $\mathbb{R}$ and suppose $F\in\mathcal{M}_n(\mu)$ and
$F\in\mathcal{M}_n(\nu)$ for some reals $\mu,\nu$. Fix any $K>0$. By Lemma~\ref{lem:key},
applied first to the center $\mu$ and then to the center $\nu$ — the quantities
$M_K(F),A_K(F),B_K(F)$, and hence $D_K(F)$, depend only on the marginal $F$ and on $K$, not
on the chosen center or coupling — we obtain
\[
\mu=D_K(F)=\nu .
\]
Therefore $\mu=\nu$. Equivalently, $\mathcal{M}_n(\mu)\cap\mathcal{M}_n(\nu)=\varnothing$
whenever $\mu\ne\nu$, which is the claimed uniqueness.
\end{proof}

\begin{proof}[Proof of Theorem~\ref{thm:main} (vector case)]
Let $F$ be a distribution on $\mathbb{R}^d$ and suppose $F\in\mathcal{M}_n(\mu)$ and
$F\in\mathcal{M}_n(\nu)$ with $\mu,\nu\in\mathbb{R}^d$. Write $\mu=(\mu^{(1)},\dots,\mu^{(d)})$
and $\nu=(\nu^{(1)},\dots,\nu^{(d)})$. Fix a coordinate $\ell\in\{1,\dots,d\}$. Let
$\pi_\ell:\mathbb{R}^d\to\mathbb{R}$ be the projection onto the $\ell$-th coordinate, and let
$F^{(\ell)}:=F\circ\pi_\ell^{-1}$ be the $\ell$-th marginal of $F$.

From $F\in\mathcal{M}_n(\mu)$ there are $X_1,\dots,X_n\sim F$ with
$\sum_{i=1}^n X_i=n\mu$ a.s. Projecting to the $\ell$-th coordinate, the scalar random
variables $X_i^{(\ell)}:=\pi_\ell(X_i)$ each have law $F^{(\ell)}$ and satisfy
$\sum_{i=1}^n X_i^{(\ell)}=n\mu^{(\ell)}$ a.s., so $F^{(\ell)}\in\mathcal{M}_n(\mu^{(\ell)})$.
Likewise $F^{(\ell)}\in\mathcal{M}_n(\nu^{(\ell)})$. By the scalar case just proved,
$\mu^{(\ell)}=\nu^{(\ell)}$. Since $\ell$ was arbitrary, $\mu=\nu$.
\end{proof}

\begin{remark}
The proof uses no integrability hypothesis on $F$. The only quantities integrated are the
\emph{bounded} clip $g_K$ and the two tail excesses $(x-K)_+$ and $(-x-K)_+$, whose finiteness
is forced by the boundedness of the total sum $\sum_i X_i\equiv n\mu$ in
\eqref{eq:cm}. Thus uniqueness of the center holds for every $n$-completely mixable
distribution, with or without a finite mean. When a mean does exist, letting $K\to\infty$ in
\eqref{eq:conservation} recovers the familiar identity $\mu=\mathbb{E}[X_1]$, since then
$A_K(F),B_K(F)\to0$ and $M_K(F)\to\mathbb{E}[X_1]$ by dominated convergence.
\end{remark}

\end{document}
